% ##############################################################################
\section{Implementation Notes}
\label{sec:implementation}

The deployed system uses AWS-managed identity and data services, a server-rendered web tier, and a modular React application. The details below are not requirements for the pattern, but they show where the sharp edges live.

\subsection{Control plane}
\label{sec:control_plane}

The control plane is an AWS Lambda API behind API Gateway. Middleware validates Cognito JWTs, resolves tenant membership from the policy store, computes effective role and module access, applies view-as guardrails, and returns the \code{/v2/me} contract. AWS Lambda Powertools provides the middleware structure and observability hooks \cite{aws_powertools}.

The endpoint uses ETag and \code{If-None-Match} semantics so repeated identity checks do not force full materialization. Full materialization still happens when the identity fingerprint changes, when the session changes, or when a page load needs a fresh contract.

\subsection{Web tier}
\label{sec:web_tier}

The web tier stores session state with iron-session \cite{iron_session}. Server components read the session, call the control plane, build CASL abilities, and render the application shell. Packed CASL rules are sent to the client for UX consistency. They are treated as hints, never authority.

\subsection{Policy store governance}
\label{sec:policy_store_governance}

Moving UI policy into DynamoDB changes the operational surface. A buggy write can alter the UI contract for an entire tenant. That is better than sprinkling tenant logic through a codebase, but it is not free.

In production, writes are IAM-scoped, reviewed, and parsed strictly at materialization. Future work should add versioned JSON Schema validation and CI checks before policy records reach production. Policy-as-data still needs governance; it simply makes the governance visible.

\subsection{Caching trade-off}
\label{sec:caching_tradeoff}

A policy edit takes effect at the next full materialization or page load. That is acceptable for tenant onboarding and module toggles, which are operator actions. For urgent security changes, push invalidation or policy-version bumps would shorten the window. We discuss this in Section~\ref{sec:futurework}.
